% Generated draft: requires mathematical and source review.
\documentclass[11pt]{article}
\usepackage[margin=25mm]{geometry}
\usepackage{fontspec}
\setmainfont{Latin Modern Roman}
\usepackage{amsmath,amssymb,mathtools,unicode-math}
\setmathfont{Latin Modern Math}
\usepackage{enumitem,needspace,xurl,hyperref}
\setlength{\parindent}{0pt}
\setlength{\parskip}{5pt}
\setlength{\emergencystretch}{3em}
\newcommand{\N}{\mathbb{N}}
\newcommand{\Z}{\mathbb{Z}}
\newcommand{\Q}{\mathbb{Q}}
\newcommand{\R}{\mathbb{R}}
\newcommand{\C}{\mathbb{C}}
\newcommand{\F}{\mathbb{F}}
\newcommand{\A}{\mathbb{A}}
\newcommand{\PP}{\mathbb{P}}
\newcommand{\E}{\mathbb{E}}
\newcommand{\eps}{\varepsilon}
\newcommand{\dd}{\,\mathrm{d}}
\begin{document}
\section*{UnsolvedMath ID: 30001341}
\textbf{Source problem code:} catalogue entry
\subsection{Definitions and mathematical statement}
Let \(V\) denote the interaction potential referred to in the supplied source. A \emph{steady state} is a time-independent solution of the underlying evolution equation; that equation is not specified in the record.

\emph{Source question.} What conditions on \(V\) permit a steady state with one peak? What conditions on \(V\) permit a steady state with two equal-sized peaks?

These are requests for existence conditions for the respective configurations. The record does not specify the evolution equation, spatial domain, boundary conditions, admissible potentials, solution class, or normalization. It also does not define a peak or whether equal size means equal height, equal mass, or another quantity. Consequently, neither the admissible steady states nor the peak conditions can be expressed as precise mathematical predicates from the supplied information. No assertion that the two configurations are mutually exclusive is supplied. Although the background discusses stability, the question does not specify a stability requirement or a notion of stability. These qualifications remain unresolved in this statement-only draft.
\subsection{Short English statement}
The question asks which interaction potentials permit a steady state with one peak or with two equal-sized peaks. The governing equation and the meaning of peak size must be supplied before this can become a fully precise mathematical problem.
\subsection{Sources}
\begin{itemize}
\item Supplied local problem registry record, permanent UnsolvedMath ID \(30001341\), source code \texttt{catalogue entry}, title \emph{Interaction Potentials for Single and Double Peak States}, in the Oberwolfach Reports Open Problems collection. No individual proposer is identified; the supplied proposed year is 2009.
\item Supplied report citation: \emph{PDE and Materials (2010)}, Oberwolfach Reports, pp.\ 2335--2404, DOI \texttt{10.4171/owr/2009/42}. This is also the literature source labelled ``original Oberwolfach report'': \url{https://doi.org/10.4171/owr/2009/42}. The supplied background describes Oberwolfach Reports as documenting mathematical workshops and their open-problem discussions.
\item Supplied literature source, labelled ``stability of peak solutions'': \url{https://www.maths.tcd.ie/EMIS/journals/EJDE/Volumes/2012/157/geigant.pdf}.
\item Supplied literature source, labelled ``MFO report text'': \url{https://ems.press/content/serial-article-files/46245?nt=1}.
\item Literature assessment dated exactly \texttt{2026-08-21}: the supplied curation labels the problem ``partially solved'' and reports existence and stability criteria and counterexamples for several interaction kernels, while reporting that a general characterization for one versus two equal peaks remains open. This is supplied dated curation, not an independently verified literature assessment; its time-sensitive status and cited claims require independent verification.
\item The supplied statement status is \texttt{corrected\_verified}. Its verification note records only that the variable denoting the interaction potential was placed in a math delimiter without changing the question; it does not supply the missing mathematical model.
\end{itemize}
\end{document}
